For every \(\mathcal S\) satisfying every premise of Wu--Bai--Cui Proposition 3.1 and Theorem 3.1, including their consistency, sequential, cross-world, proxy-exclusion, treatment and mediator positivity, observed-bridge-existence, and completeness conditions, put \(P=P_{\mathcal S}^O\). If \(\mathcal H_{\mathrm{obs}}(P)\neq\varnothing\), Wu--Bai--Cui Proposition 3.1 makes every \(x\in\mathcal H_{\mathrm{obs}}(P)\) latent-valid for \(\mathcal S\), and their Theorem 3.1 gives \(\psi(\mathcal S)=\ell(P)^{\top}x=\mathbb E[h_1(W)\mid A=0]\). Hence \(\ell(P)^{\top}x\) is constant over \(\mathcal H_{\mathrm{obs}}(P)\), \(\mathsf{RI}(P)\) holds, and \(\psi(\mathcal S)=r(P)^{\top}b(P)\), whether or not \(B(P)\) is nonsingular. When \(\mathcal S\in\mathfrak M_{\mathrm{WBC}}^{+}\) also holds, this is exactly the conditional-envelope conclusion.